\section{Extended related work}
\label{app:related}

This appendix records background that informed the design but that the main
argument does not depend on. Nothing here supports a result in
\S\ref{sec:results}; a reader who wants only the load-bearing literature should
read \S\ref{sec:related} and stop there. The organising convention is that each
entry says what the work does and then says which of three things this study
did with it: \textbf{uses} it as method, \textbf{cites} it as context, or
\textbf{declines} it. \S\ref{app:rel-declined} collects the declined set,
because in a study whose thesis is that the binding constraint is information
rather than machinery, the machinery left out is part of the argument.

\subsection{Further informal sources on the game}
\label{app:rel-informal}

Beyond Pagat \cite{pagat} and Develin \cite{develin} the descriptive corpus is
thin. Pagat's tactics list covers deliberate misses that inform a partner,
locking out a dangerous opponent by never asking them anything, and exhausting
one rank across all opponents before switching; it also records a
partner-signalling convention and a ``Forced Claims'' rule, both attributed to
named contributors. Develin's chapter contains the corpus's only quantitative
endgame discussion, in which the alternative to a coin-flip declaration would
need a success probability greater than one. The Wikipedia entry
\cite{wikiliterature} names the ``stalemate breaker'', and secondary strategy
pages \cite{dorsa,strategy-site} restate blackballing, asking the asker, and an
endgame delegation heuristic. These are used as evidence about how the game is
played and talked about, not as sources of quantitative claims. The two
substantial sources also disagree with each other on the point that matters most
to a study of conventions: Develin's chapter forbids pre-agreed conventions
outright, where Pagat records a partner-signalling convention as an accepted
tactic. That disagreement is why the partner regime is a declared parameter in
\S\ref{sec:mech-partner} rather than a fixed assumption.

On Somani's agent \cite{somani,somani-server}: it is a Python implementation
with a complete rules engine and a sound deduction fixed point, and the
four-player, 48-card variant it implements makes the declaration problem
structurally easier than the one studied here. With a single teammate, naming
the allocation of a six-card half-suit ranges over $2^6$ possibilities and is
frequently forced by the deduction fixed point; with two teammates it ranges
over $3^6 = 729$, and the allocation question stops being a formality. No
strength numbers have been published for it, so it is not usable as a baseline
in any case. The v0.4 report records a reading of its training code that
suggests the released model never received a win/loss gradient; that reading is
repeated here without re-verification, because nothing in this study depends on
it.

Valet \cite{goadrich2026valet} is the one piece of 2026 infrastructure that
could change the situation. It encodes twenty-one traditional
imperfect-information card games in a common description language and publishes
measured branching factors and game durations, which is exactly the external
calibration a claim like ``this game is hard'' needs. Whether any of the
twenty-one is a team game is not stated on its abstract page, and the full text
was not fetched, so no comparison is drawn anywhere in this report. Recorded so
that the next study fetches it rather than re-deriving the question.

\subsection{Determinization and perfect-information sampling}
\label{app:rel-determinization}

\S\ref{sec:related-determinization} states that averaging perfect-information
values over deals sampled from the public history collapses to a hit-probability
heuristic in this game. The collapse is the following. Almost every reachable
position carries the same perfect-information value, because the team on turn
takes everything it can reach, so with $t$ ranging over opponents, $c$ over
askable cards and $h$ the public history,
\[
  \arg\max_{(t,c)} \; \mathbb{E}\bigl[V^{\mathrm{PI}}(t,c)\bigr]
  \;\approx\;
  \arg\max_{(t,c)} \; \Pr[\,t \text{ holds } c \mid h\,].
\]
The right-hand side is one ply deep and blind to the information an ask leaks,
the information a refusal supplies and the option value of postponing a
declaration. This is a degeneracy of the perfect-information game itself, and it
is why the determinized search of \S\ref{sec:search} scores candidates by
rolling out under a belief-limited blueprint rather than by solving the sampled
world.

Determinization for bridge was proposed well before it was made to work
\cite{levy}, and constrained deal generation in practice is done by rejection:
deal against restrictions, then discard deals the program's own bidding module
would not have produced \cite{ginsberg-gib}. Frank and Basin's diagnosis
--- strategy fusion and non-locality --- is what makes it unsound however many
worlds are drawn \cite{frank-basin-ai}, and the repairs that work are those
attacking non-locality as well as fusion \cite{frank-basin-aaai}. Long,
Sturtevant, Buro and Furtak's leaf correlation, bias and disambiguation factor
explain when it nevertheless performs well \cite{long-pimc}. Information Set
Monte Carlo Tree Search replaces the independent trees with a single tree over
information sets, redeterminizing each iteration \cite{cowling-ismcts}, and has
been extended by re-determinizing at non-root seats \cite{goodman-ris}.
\textbf{Cited as context}; none is used as an engine here.

$\alpha\mu$ \cite{cazenave2019alphamu} deserves a separate note because it is
the closest thing in this literature to a repair this study could have adopted.
It backs up Pareto fronts of boolean world-vectors instead of scalar averages,
which represents ``this plan works in world set A, that plan works in world set
B'' without committing to a scalar mean --- structurally the same move as the
multi-valued leaf of \S\ref{sec:mech-leaves}, one ply down --- and its published
timings are an existence proof that vector-valued search over roughly twenty
worlds is affordable in about a second per move on a CPU. Its reported gain in
Bridge comes from differing with plain sampling on only a small fraction of
decisions, which is the shape a good search should have and the calibration
\S\ref{sec:search} uses when judging whether its own search is mis-scaled.
\textbf{Cited as context, and its diagnostic used}: the algorithm is declined
because it repairs unsoundness rather than degeneracy, and this game suffers the
latter.

Knowledge-based paranoia search \cite{edelkamp2021paranoia} is the closest
methodological neighbour this study has: it runs on a CPU with no network, is
built around deduced knowledge rather than sampled worlds, and reasons about the
worst case over consistent worlds rather than the average. All three match this
engine, whose deduction state carries hard exclusions, capacity constraints and
ask-legality certificates, and whose count law already computes the probability
of a named allocation. Worst-case-over-consistent-worlds is also the right
criterion for the one irreversible action in the game. \textbf{Its criterion is
used} in the declaration path of \S\ref{sec:mech-declare}; its game-tree search
is declined, because Skat's horizon and branching factor are not this game's. No
compute cost is attributed to it anywhere, because its runtime is not stated on
its abstract page.

\subsection{Depth-limited and subgame solving}
\label{app:rel-depthlimited}

Brown, Sandholm and Amos \cite{brown2018depthlimited} is the one result in this
group that is \textbf{used as method}. Its content is a single observation with
a large consequence: a search that stops at a depth limit and evaluates the leaf
with one value function has silently assumed the opponent plays one fixed
continuation, and that assumption is what makes depth-limited search in an
imperfect-information game unsound. Letting the opponent choose among $k$
continuations at the limit, each yielding its own leaf values, forces robustness
to all of them, and the demonstration ran on a four-core CPU and 16\,GB rather
than a supercomputer. The prerequisite it identifies as expensive --- a
portfolio of distinct continuation strategies --- is already paid for in this
project, because a set of compiled baseline policies of different styles exists
and is used for evaluation. The caveat is stated in \S\ref{sec:mech-leaves} and
repeated here: the soundness argument is two-player zero-sum, this game's leaf
has three uncorrelated opponent seats, and the construction is therefore adopted
as a robustification heuristic. A second caveat is that the method assumes a
blueprint approximating equilibrium, and this policy is a fitted heuristic, so
what the leaves buy is robustness against a style portfolio and not robustness
against equilibrium deviations.

Zhang and Sandholm's knowledge-limited subgame solving
\cite{zhang2021subgamesolvingwithoutck} is \textbf{cited as context and its
warning honoured}. Their motivation is precisely this game's blocker: current
techniques analyse the entire common-knowledge closure, which is acceptable when
that closure is small and impossible when it is not. Their remedy is to work
with low-order knowledge, pruning nodes no longer reachable on arrival at an
information set, and their honest result is that this can increase
exploitability, after which they develop the conditions that restore safety.
This engine already ships a crude instance of the same move --- a gate that
restricts candidate asks to those with strictly positive hard-consistent
probability --- whose measured gains were large and whose exploitability cost
has never been measured. \S\ref{sec:mech-exploit} exists because of this
paragraph.

Subgame solving in adversarial team games \cite{zhang2022teamsubgame} builds a
gadget from the team belief DAG and refines a blueprint by column generation,
running in polynomial time given a best-response oracle when the blueprint is
sparse. \textbf{Declined}, and recorded as a verified non-transfer: the sparsity
condition is the escape hatch, and a policy with support over roughly seventy
legal asks at each of roughly a hundred decisions per game is not sparse. It is
named explicitly because it is the most natural thing a reader will ask about.

The counter-strategy line is \textbf{cited for one warning}. Continual
depth-limited responses \cite{milec2021cdr} state plainly that existing robust
response methods do not compose with limited-look-ahead solving off the shelf,
and supply a fix; adapting beyond the depth limit \cite{milec2025abd} extends
the idea to counter-strategies in large games while staying robust against
rational opponents. Any future version of this policy that carries both a
depth-limited search and an online opponent model needs that composition
argument, and this report does not have one, because it builds only the first.
Look-ahead reasoning with a learned model \cite{kubicek2025lamir} is the right
direction and requires a learned abstraction, so it is deferred. Equilibrium
refinements in gadget games \cite{kubicek2026refinements} are not applicable
here at all --- no gadget game is solved anywhere in this study --- but the
finding transfers as a meta-lesson: when a subproblem has many optima that are
theoretically equivalent, the choice among them is worth more than half of the
exploitability, and that is a strong statement of the tie-breaking argument in
\S\ref{sec:diag-ties}.

The full-solver line is \textbf{cited as context and declined as
implementation}. DeepStack \cite{moravcik2017deepstack} introduced continual
re-solving from a range and a set of opponent counterfactual values; its
re-solve gadget requires one value per opponent information set, which at three
opponents holding nine cards each is not reachable on any CPU budget. ReBeL
\cite{brown2020rebel} treats the public belief state as the state and proves
that running the same procedure at test time is safe, which is the licence under
which \S\ref{sec:search} operates at all; its algorithm iterates over the belief
support and is therefore blocked by cardinality. Student of Games
\cite{schmid2023sog} grows the public tree during solving and identifies
Scotland Yard --- a public-observation, hidden-position game --- as the closest
published analogue to this information structure; the transferable piece is
architectural and cheap, namely that a candidate set should be grown where the
search is uncertain rather than fixed at a constant, and even that is not
implemented here.

\subsection{Policy-aware inference in trick-taking games}
\label{app:rel-inference}

The Skat line is the direct ancestor of this study's belief construction and is
\textbf{used}: reweighting the deals consistent with the public record by a
model of how each player would have acted, whether through a supervised
per-card location model \cite{solinas-supervised} or through reach probabilities
taken from a policy \cite{rebstock-inference}, and reported as worth $+217$
tournament points per 36 games as a standalone module \cite{buro-kermit}. The
policy prior of \S\ref{sec:belief-prior} is a two-parameter version of the same
idea. What \S\ref{sec:belief-exactworse} contributes to this line is a case
where the trade-off inverts far enough to be visible in win rate: exactness is
affordable here, and the exact object still loses, because it is the one that
cannot see behaviour.

History filtering --- constructing a history consistent with a public state ---
is FNP-complete in the worst case, and enumeration is polynomial only for sparse
public trees \cite{solinas-history}. \textbf{Cited as context}, and this game
escapes the problem entirely: its public record determines every card movement,
so the only object to be reconstructed is the initial deal, which is what makes
the exact count law of \S\ref{sec:belief-count} possible.

Belief modelling in cooperative games is \textbf{declined}. The Bayesian Action
Decoder made the public belief an explicit object of the policy
\cite{foerster2019bad}, and its independent-marginals factorisation is the part
that does not survive contact with fixed hand sizes and card conservation; the
repair is the constrained scaling and block enumeration this engine already
implements. Approximate policy iteration over common-payoff games
\cite{sokota2021capi} formalises factorised prescriptions and is a formalism
rather than a tool here. Learned Belief Search \cite{hu2021lbs} learns an
approximate belief because Hanabi's exact belief is intractable; this game's is
not, so learning one would be strictly worse than what exists. That is a real
structural advantage of this domain and is stated plainly in
\S\ref{sec:related-inference}: the belief half of the problem is solved, and
only the search half is open.

\subsection{Cooperative play, conventions, and deviations}
\label{app:rel-cooperative}

Hanabi is the reference benchmark for coordination through observable actions
\cite{bard-hanabi}, and is relevant because this game likewise has no side
channel. The Simplified Action Decoder showed that letting teammates condition
on the greedy rather than the exploratory action is a prerequisite for
conventions to survive exploratory training \cite{hu-sad}; Other-Play
demonstrated how much of a self-play score is convention rather than skill, with
the same agents scoring far worse when paired across independent training runs
\cite{hu-otherplay}; and Off-Belief Learning provides an explicit dial on
convention depth together with a uniqueness result that makes independent runs
interoperable \cite{hu2021obl}. Test-time search over a fixed blueprint provides
a further improvement without any training at all \cite{lerer2020sparta}. Joint
policy search, which scores simultaneous changes at several information sets ---
which is what inventing a convention requires --- improved a bridge bidding
system by $+0.63$ IMPs per board against a strong commercial program
\cite{tian2020jps}. All \textbf{cited as context}; none is trained here.

Two structural asymmetries separate this game from Hanabi and are worth stating
because they cut in opposite directions. The first is the audience. Hanabi has
no adversary, so every bit of signal is pure profit; here a leaked card location
is directly executable, because once it is public that a player holds a card,
any opponent holding another card of that half-suit can take it and keep the
turn. The Hanabi construction that most clearly fails to port is the modular
hat-guessing code, which works because each receiver sees every other receiver's
hand and can subtract it \cite{cox-hanabi}; here no player sees any hand. The
second asymmetry runs the other way. Off-Belief Learning's grounded level-one
play, which refuses to read conventional meaning into a partner's action, is an
approximation in Hanabi and is nearly exact here, because ask legality is a
rules-level certificate: a legal ask proves its asker holds another card of the
named half-suit, with no convention involved. The prediction that grounded play
is much closer to optimal in this game than in Hanabi is testable with a single
flag and no training, and it is the ablation of \S\ref{sec:results-ablations}
that separates hard certificate reading from the soft prior weights.

Self-explaining deviations \cite{hu2022sed} are the most on-topic piece of this
literature for the question of finding non-obvious plans, and are
\textbf{declined for this version with the reason recorded}. A self-explaining
deviation is an action that departs from what common understanding calls
reasonable, taken so that a teammate reasoning by theory of mind concludes that
circumstances must be abnormal; the associated algorithm produced the first
algorithmic finesse plays in Hanabi. The analogue here is exact and currently
unexploited: an ask certifies a holding, so an ask with low hit probability that
certifies something a teammate could not otherwise deduce is precisely such a
deviation, and the belief machinery to score one already exists. It is declined
in this study for two reasons. The information features of the shipped policy
are all negative --- it charges for information rather than paying for it --- so
the mechanism would have to be built rather than tuned; and a deviation of this
kind is a convention by another name, which means it is the mechanism most
likely to look excellent in self-play and fail against a human or an earlier
version. Any future attempt must be reported with its cross-play cell and gated
behind the partner regime of \S\ref{sec:mech-partner}.

\subsection{Human-compatible play and regularised search}
\label{app:rel-human}

piKL \cite{jacob2022pikl} regularises a search result toward a reference policy
with a divergence bound that shrinks as the regularisation weight grows, and
reports improved top-one agreement with human moves in chess and Go alongside
strong play. \textbf{Cited as context, and its shape borrowed.} The framing that
matters here is not human-likeness: it is that a KL anchor is a tunable,
measurable interpolation between a reference policy and an unconstrained search,
with the reference recovered exactly in the limit, so an ablation along that
axis has a guaranteed floor. In a team game the anchor also does real strategic
work, because a search that finds a clever line its own teammates will misread
is a net loss.

Human-regularised search and learning \cite{hu2022humanregularized} extends this
to coordination with real partners in three steps, of which the third ---
regularised search at test time to absorb partner distribution shift --- is the
published architecture for exactly the requirement that this policy behave
differently with a bot partner and a human one. \textbf{Cited as context};
there is no human corpus for this game, so it is not implemented. Cicero
\cite{fair2022cicero} is cited for the family rather than the architecture: a
seven-player mixed cooperative-competitive game whose planning core is an
anchored search over actions, and whose lesson is that in games with allies the
planner's job is to act well against a model of what the other players will
actually do rather than against an equilibrium. The ad-hoc human-AI coordination
challenge \cite{dizdarevic2025ah2ac2} supplies the negative instruction this
study follows: an agent trained with no human data at all outperformed a
best-response-to-behavioural-cloning baseline against human proxies, so a small
human corpus is not the way in.

\subsection{Team games and common information}
\label{app:rel-team}

The common-information approach \cite{nayyar2013common} reformulates a
decentralised control problem as a single-agent problem for a coordinator who
knows the common information and selects prescriptions --- maps from each
controller's local information to its action --- and shows that the resulting
problem is a POMDP. \textbf{Used as a correctness lens.} Two consequences are
load-bearing in this report. First, it names the design object: the public
deduction state is one shared thing, and what distinguishes seats is a
refinement of it, which is the mechanism of \S\ref{sec:mech-knowledge} and the
reason that mechanism is $O(1)$ rather than six-fold. Second, it explains why
teammates must run identical procedures: a prescription is agreed in advance,
and a seat that deviates by searching invalidates its teammates' belief updates.
The coordinator's own problem is not solvable at this scale --- a prescription
maps every nine-card hand to an ask --- so nothing here computes one.

The equilibrium-computation literature for teams is \textbf{cited to delimit
scope and declined in full}. Celli and Gatti establish the complexity of the
team-maxmin equilibrium with correlation and the worst-case cost of forgoing
correlation \cite{celli-gatti}; Farina and coauthors connect ex-ante team
collusion to extensive-form correlation \cite{farina-collusion}; the team belief
DAG and the two-player conversion give a zero-sum game whose equilibria are
realization-equivalent to a team-maxmin equilibrium with correlation
\cite{zhang-tbdag,carminati-tpi}, at a representation cost exponential in how
much private information distinguishes teammates inside a public state. That
exponent is the blocker: three teammates hold unrestricted nine-card hands, and
one seat's information set at deal time contains
$45!/(9!)^5 \approx 1.9 \times 10^{28}$ deals. Hidden-role games
\cite{carminati2024hiddenrole} are cited only to delimit scope: teams here are
public from the deal, so there is no role uncertainty and none of that machinery
applies. Nothing in this report is an equilibrium result, and no claim in it
should be read as one.

\subsection{Counting and constrained sampling}
\label{app:rel-counting}

Counting the deals consistent with a public history is a permanent computation.
Ryser's inclusion--exclusion formula is exact in $O(2^n n)$ time \cite{ryser},
which is unusable at $n = 54$, and Glynn's sign-vector formula is the other
exact method \cite{glynn}; the Jerrum--Sinclair--Vigoda chain is a fully
polynomial randomised approximation scheme \cite{jsv} whose constants make it
impractical at this scale. Matrix scaling is the usual practical approximation
and carries a deterministic strongly polynomial guarantee \cite{lsw}; a Sinkhorn
iteration of this kind is what the deployed approximate path uses. Sequential
importance sampling is the standard contingency-table alternative
\cite{chen-sis}, and is not used, because it can underestimate badly while
appearing to have converged. \textbf{Cited as context, with one used.}

The reason none of this is on the critical path is structural. The hidden state
is exactly the initial deal, each seat's share of it is a known constant number
of cards, and every ask-legality certificate is confined to a single half-suit,
so the counting problem factorises into blocks that can be enumerated exactly
within a half-suit and combined by a capacity dynamic programme. That is the
construction of \S\ref{sec:belief-count}, and it is why this study can ask
whether exactness helps rather than whether it is affordable.

\subsection{Evaluation methodology and variance reduction}
\label{app:rel-eval}

The v0.4 study declined variance reduction outright. This study takes the same
decision about fitted control variates and a different one about pairing, and
the distinction is worth stating because the two are often filed together.
Advantage-sum estimators \cite{zinkevich-divat} and AIVAT \cite{burch-aivat}
learn or hand-build a value function and subtract it, and their reported gains
are large in two-player settings and smaller in six-player ones, which is the
closer analogue here. They remain \textbf{declined}. AIVAT is unbiased for
\emph{any} value function, which is a constraint on its expectation rather than
on any realised estimate, and a control variate fitted by the same author on the
same agent using the same evaluation data is the pathology that literature
itself warns about.

What this study adopts instead is pairing on common random numbers: the same
deals played in both arms, with treatment and control alternated across seat
rotations, and the difference taken per deal
(\S\ref{sec:protocol-paired}, \S\ref{sec:fit-paired}). It costs nothing to
implement, introduces no fitted object, and removes the largest part of deal
variance --- the report quotes the realised width ratio against the unpaired
estimator rather than asserting a benefit. It is also the standard remedy rather
than a novel one: duplicate bridge replays each board with the same cards in the
same seats, and the computer poker competition added common seeds across
pairings for the same reason \cite{acpc}. The cost of the alternative is that a
fitted estimator would have to be validated, and validating it would consume
more games than pairing saves.

One further piece of methodology is named because it is frequently applied here
by analogy and should not be. Refining an abstraction does not monotonically
reduce exploitability \cite{waugh-abstraction}, which sounds like the general
form of \S\ref{sec:belief-exactworse}'s finding that substituting a more
accurate belief into a policy fitted around a less accurate one does not improve
it. The analogy is suggestive and the result does not transfer: this policy is
not an abstraction of a solved game, so neither the pathologies that literature
describes nor its guarantees apply. The finding here is measured on its own
terms and defended by a control, not inherited.

Rating and stopping are treated briefly. Bradley--Terry-style rating models
\cite{coulom-whr,herbrich-trueskill} are not invariant to redundancy in the
rated population and cannot represent cyclic dominance, both of which are real
failure modes when the rated
population is one training run's own checkpoints; maximum-entropy Nash averaging
over the antisymmetric logit matrix is the proposed alternative
\cite{balduzzi-reeval}. Ratings appear in this report only as a compact summary
of a fixed round-robin, never as a claim about a strategy space. And monitoring
an interval computed for a fixed sample size and stopping when it looks
favourable is not a test, which is why sequential designs on normalised Elo
exist \cite{nelo} and why the seed banks here are fixed in advance and reported
at their planned size (\S\ref{sec:protocol-seeds}).

\subsection{Work this study deliberately does not use}
\label{app:rel-declined}

The declined set, with the reason for each. These are design decisions, and
stating them is the point of this appendix.

\begin{itemize}
  \item \textbf{Full solvers over public belief states}
        \cite{moravcik2017deepstack,brown2020rebel,schmid2023sog}. Every one
        requires iterating or learning over a common-knowledge closure whose
        cardinality here is on the order of $10^{28}$. The closure is used as a
        sufficient statistic for evaluation and never as an index set for
        iteration. ReBeL's safety result is retained as the licence for
        test-time search; its algorithm is not.
  \item \textbf{The team belief DAG and subgame solving over it}
        \cite{zhang-tbdag,carminati-tpi,zhang2022teamsubgame}. Representation
        cost is exponential in the private information distinguishing teammates
        inside a public state, and the polynomial escape hatch is conditional on
        blueprint sparsity that this policy does not have.
  \item \textbf{Learned belief models} \cite{hu2021lbs,foerster2019bad}. The
        exact belief is computed and oracle-validated here, so an approximation
        would be strictly worse. This is the one place where this domain is
        easier than Hanabi rather than harder.
  \item \textbf{Neural policies and decision-time gradient updates}
        \cite{fickinger2021rlfinetuning}. There is no network training
        infrastructure in this project and none is warranted by the results:
        the two largest effects measured in this report are a tie-break and a
        statistical guard, neither of which is a capacity problem. The 2026
        lightweight-agent study reaches the same conclusion independently
        \cite{kelidari2026goldstandard}.
  \item \textbf{Hidden-role machinery} \cite{carminati2024hiddenrole}. Teams are
        public from the deal; there is no role uncertainty to model.
  \item \textbf{Gadget-game equilibrium refinement}
        \cite{kubicek2026refinements} and \textbf{learned-model look-ahead}
        \cite{kubicek2025lamir}. No gadget game is constructed here and no model
        is learned; both are recorded for their transferable lessons only.
  \item \textbf{Fitted control variates for evaluation}
        \cite{zinkevich-divat,burch-aivat}. Declined for the reason in
        \S\ref{app:rel-eval}; pairing on common random numbers is adopted
        instead.
  \item \textbf{Self-explaining deviations} \cite{hu2022sed}. Declined for this
        version rather than permanently, with the construction sketched in
        \S\ref{app:rel-cooperative} so that the next study can build it against
        a cross-play measurement rather than a self-play one.
\end{itemize}

Two exclusions are not in this list because they are not choices. There is no
equilibrium computation anywhere in this report, and no claim in it is an
equilibrium claim. And there is no human dataset, so every statement about human
play in \S\ref{app:human} is a statement about a small number of recorded games
and is scoped as such.
